\documentclass{amsart}
% For dealing with references we use the comment environment
\usepackage{verbatim}
\newenvironment{reference}{\comment}{\endcomment}
%\newenvironment{reference}{}{}
\newenvironment{slogan}{\comment}{\endcomment}
\newenvironment{history}{\comment}{\endcomment}

% For commutative diagrams we use Xy-pic
\usepackage[all]{xy}

% We use 2cell for 2-commutative diagrams.
\xyoption{2cell}
\UseAllTwocells

% We use multicol for the list of chapters between chapters
\usepackage{multicol}

% This is generally recommended for better output
\usepackage{lmodern}
\usepackage[T1]{fontenc}

% For cross-file-references

% Package for hypertext links:
\usepackage{hyperref}

% For any local file, say "hello.tex" you want to link to please

% Theorem environments.
%
\theoremstyle{plain}
\newtheorem{theorem}[subsection]{Theorem}
\newtheorem{proposition}[subsection]{Proposition}
\newtheorem{lemma}[subsection]{Lemma}

\theoremstyle{definition}
\newtheorem{definition}[subsection]{Definition}
\newtheorem{example}[subsection]{Example}
\newtheorem{exercise}[subsection]{Exercise}
\newtheorem{situation}[subsection]{Situation}

\theoremstyle{remark}
\newtheorem{remark}[subsection]{Remark}
\newtheorem{remarks}[subsection]{Remarks}

\numberwithin{equation}{subsection}

% Macros
%
\def\lim{\mathop{\mathrm{lim}}\nolimits}
\def\colim{\mathop{\mathrm{colim}}\nolimits}
\def\Spec{\mathop{\mathrm{Spec}}}
\def\Hom{\mathop{\mathrm{Hom}}\nolimits}
\def\Ext{\mathop{\mathrm{Ext}}\nolimits}
\def\SheafHom{\mathop{\mathcal{H}\!\mathit{om}}\nolimits}
\def\SheafExt{\mathop{\mathcal{E}\!\mathit{xt}}\nolimits}
\def\Sch{\mathit{Sch}}
\def\Mor{\mathop{\mathrm{Mor}}\nolimits}
\def\Ob{\mathop{\mathrm{Ob}}\nolimits}
\def\Sh{\mathop{\mathit{Sh}}\nolimits}
\def\NL{\mathop{N\!L}\nolimits}
\def\CH{\mathop{\mathrm{CH}}\nolimits}
\def\proetale{{pro\text{-}\acute{e}tale}}
\def\etale{{\acute{e}tale}}
\def\QCoh{\mathit{QCoh}}
\def\Ker{\mathop{\mathrm{Ker}}}
\def\Im{\mathop{\mathrm{Im}}}
\def\Coker{\mathop{\mathrm{Coker}}}
\def\Coim{\mathop{\mathrm{Coim}}}

% Boxtimes
%
\DeclareMathSymbol{\boxtimes}{\mathbin}{AMSa}{"02}

%
% Macros for moduli stacks/spaces
%
\def\QCohstack{\mathcal{QC}\!\mathit{oh}}
\def\Cohstack{\mathcal{C}\!\mathit{oh}}
\def\Spacesstack{\mathcal{S}\!\mathit{paces}}
\def\Quotfunctor{\mathrm{Quot}}
\def\Hilbfunctor{\mathrm{Hilb}}
\def\Curvesstack{\mathcal{C}\!\mathit{urves}}
\def\Polarizedstack{\mathcal{P}\!\mathit{olarized}}
\def\Complexesstack{\mathcal{C}\!\mathit{omplexes}}
% \Pic is the operator that assigns to X its picard group, usage \Pic(X)
% \Picardstack_{X/B} denotes the Picard stack of X over B
% \Picardfunctor_{X/B} denotes the Picard functor of X over B
\def\Pic{\mathop{\mathrm{Pic}}\nolimits}
\def\Picardstack{\mathcal{P}\!\mathit{ic}}
\def\Picardfunctor{\mathrm{Pic}}
\def\Deformationcategory{\mathcal{D}\!\mathit{ef}}

\setlength{\emergencystretch}{3em}
\begin{document}
\title[Stacks Project, Tag 0FIW]{267 --- Equidimensionality of catenary Noetherian local S2 rings}
\maketitle
\noindent Copyright (C) 2005--2025 Johan de Jong. Stacks Project authors. Source: \href{https://stacks.math.columbia.edu/tag/0FIW}{Tag 0FIW}.

\noindent Editorial difficulty note (not author text). Difficulty H2: The proof separates maximal-dimensional from lower-dimensional irreducible components and chooses a generic intersection point. Catenarity forces a local dimension at least two, while S2 forces depth at least two; the resulting punctured spectrum is disconnected, contradicting depth-two connectedness. The geometric contradiction is more than a routine dimension formula..

\noindent Editorial provenance note (not author text). History: excerpt from revision \texttt{a0444\allowbreak 6e57e\allowbreak c1fbc\allowbreak 252a8\allowbreak 71afc\allowbreak ec775\allowbreak 2fb28\allowbreak 07b14}, local-cohomology.tex, lines 353--384. Reference presentation was adapted; original-author context and core support blocks were retained or added. The mathematical wording is unchanged. Licensed under GNU FDL 1.2 or later; no invariant sections or cover texts. See ../licenses/COPYING. Context blocks reproduce author definitions and supporting results; they are not additional counted problems.

\begin{lemma}
\label{lemma-catenary-S2-equidimensional}
\begin{reference}
\href{https://stacks.math.columbia.edu/bibliography}{Bibliography: EGA (IV Corollary 5.10.9)}
\end{reference}
Let $A$ be a Noetherian local ring which is catenary and $(S_2)$.
Then $\Spec(A)$ is equidimensional.
\end{lemma}

\begin{proof}
Set $X = \Spec(A)$. Say $d = \dim(A) = \dim(X)$. Inside $X$ consider the
union $X_1$ of the irreducible components of dimension $d$ and the union
$X_2$ of the irreducible components of dimension $< d$. Of course
$X = X_1 \cup X_2$. If $X_2 = \emptyset$,
then the lemma holds. If not, then $Z = X_1 \cap X_2$ is a nonempty closed
subset of $X$ because it contains at least the closed point of $X$.
Hence we can choose a generic point $z \in Z$ of an irreducible component
of $Z$. Recall that the spectrum of $\mathcal{O}_{Z, z}$ is the set of points
of $X$ specializing to $z$. Since $z$ is both contained in an
irreducible component of dimension $d$ and in an irreducible component
of dimension $< d$ we obtain nontrivial specializations $x_1 \leadsto z$ and
$x_2 \leadsto z$ such that the closures of $x_1$ and $x_2$ have different
dimensions. Since $X$ is catenary, this can only happen if at least
one of the specializations $x_1 \leadsto z$ and $x_2 \leadsto z$ is not
immediate! Thus $\dim(\mathcal{O}_{Z, z}) \geq 2$. Therefore
$\text{depth}(\mathcal{O}_{Z, z}) \geq 2$ because $A$ is $(S_2)$.
However, the punctured spectrum $U$ of $\mathcal{O}_{Z, z}$ is disconnected
because the closed subsets $U \cap X_1$ and $U \cap X_2$ are disjoint
(by our choice of $z$) and cover $U$. This is a contradiction with
Lemma \href{https://stacks.math.columbia.edu/tag/0BLR}{lemma depth 2 connected punctured spectrum (Tag 0BLR)}
and the proof is complete.
\end{proof}
\end{document}
